\section{Proof of the Sharp Progressive-Fault Limit}
\subsection{A schedule-uniform converse}
\begin{proof}[Converse in Theorem~\ref{thm:sharp}]
For any feasible healthy path chosen to be zero on the prefix,
\begin{equation}
 2\sigma^2D_H^\infty(\pi)\ge
 \sum_{j\notin I_\pi}a_j^2+
 \sum_{j\in I_\pi}(2\epsilon_j a_jz_j-z_j^2).\label{eq:adv}
\end{equation}
The inequality follows by evaluating one feasible point in the projection objective.

Consider $L$ consecutive post-onset slots with minimum amplitude $A>0$. Let $d$ slots be excluded and let the clipped observed-run lengths be $n_1,\ldots,n_m$. Then $m\le d/k+1$. On every run use the zero-endpoint signed tent $z_i=\epsilon r\min(i,n-1-i)$. All pieces join to zero across block boundaries and excluded intervals. Set
\begin{equation}
 C=\frac r4\left(2A-\frac{rL}{2}\right).
\end{equation}
When $C\ge0$, tent height is at most $rL/2$, and
$\sum_{i=0}^{n-1}\min(i,n-1-i)\ge(n^2-2n)/4$. The local contribution to \eqref{eq:adv} is at least
\begin{align}
 A^2d+C\sum_i n_i^2-2C(L-d)
 &\ge A^2d+\frac{Ck(L-d)^2}{d+k}-2CL.\label{eq:macro1}
\end{align}
The right side remains valid when there are no observations. Put $y=d+k$ and apply AM--GM to
\begin{equation}
 (A^2+Ck)y+\frac{Ck(L+k)^2}{y}
 -2Ck(L+k)-A^2k-2CL.
\end{equation}
Thus every calendar has a feasible local explanation giving at least
\begin{equation}
 \begin{split}
 \Psi(A,L)={}&2(L+k)\sqrt{Ck(A^2+Ck)}\\
 &-2Ck(L+k)-A^2k-2CL.
\end{split}\label{eq:macro}
\end{equation}
The nonnegative part may be used, or zero when $C<0$.

Fix $\epsilon\in(0,1)$ and partition the late interval $[\lceil\epsilon H\rceil,H]$ into blocks of length $L=\lfloor H^{3/4}\rfloor$. Uniformly there, $A=\Theta(H)$ and $rL/A=o(1)$, so $C=(rA/2)(1+o(1))$. The leading sum in \eqref{eq:macro} is
\begin{equation}
 \sqrt{2kr}\sum_{\rm blocks}A^{3/2}L.
\end{equation}
The accumulated $A^2k$ terms are $O(H^3/L)=O(H^{9/4})$, the $CL$ and $CkL$ terms are $O(H^2)$, and replacing $C$ by $rA/2$ incurs $O(H^{9/4})$. The last incomplete block contributes only a lower-order remainder. The Riemann sum is
\begin{equation}
 \sum A^{3/2}L
 =\tfrac25\eta^{3/2}(1-\epsilon^{5/2})H^{5/2}+o(H^{5/2}).
\end{equation}
These bounds are uniform in the calendar because only its run-count inequality was used. Taking the infimum deficit before the limit and then letting $\epsilon\downarrow0$ gives
\begin{equation}
 \liminf_{H\to\infty}
 \frac{2\sigma^2D_H^{\infty,*}}{\sqrt{kr}\eta^{3/2}H^{5/2}}
 \ge\frac{2\sqrt2}{5}.\label{eq:sharp_lower}
\end{equation}
The converse remainder used here is $o(H^{5/2})$.
\end{proof}

\subsection{Symmetric local projection and affine cancellation}
\begin{proof}[Constructive upper bound]
For an even $n=2m$, let $L=2(n+k)$ and $c=(L+1)/2$. Local admitted times are
\begin{equation}
\begin{split}
 I={}&\{1,\ldots,m\}\cup\{m+k+1,\ldots,3m+k\}\\
 &\cup\{3m+2k+1,\ldots,L\}.
\end{split}
\end{equation}
Their signs are $+,-,+$. Write $c_0=(k+1)/2$ and, in observation order, define
\begin{align}
 z^+_{i}&=r(c_0+m-i),&&1\le i\le m,\nonumber\\
 z^-_{i}&=-r[c_0+\min(i-1,n-i)],&&1\le i\le n,\nonumber\\
 z^{+,R}_{i}&=r(c_0+i-1),&&1\le i\le m.\label{eq:localz}
\end{align}
Let $a\ge r(k+n-1)/2$, $s_i^0=a\epsilon_i$, and $w=s^0-z$. Within a half-window the slope magnitude is $r$; across a transition the change is $2rc_0=r(k+1)$; the central edge of the negative window is flat. Thus $z$ is rate-feasible. Each level appears twice with each sign, so $\sum_iw_i=0$.

Let $Q_i=\sum_{j\le i}w_j$. The left-side partial sums are nonnegative, reach zero at the center of the negative window, and become nonpositive on the reflected right side. Whenever $Q_i\ne0$, the corresponding slope of $z$ equals $-r\operatorname{sign}(Q_i)$ per unit time. Summation by parts gives
\begin{equation}
 h_{\K_I(\infty,r)}(w)
 =r\sum_i(t_{i+1}-t_i)|Q_i|=w^Tz.
\end{equation}
Equivalently, $q_i=-Q_i$ satisfies $z-s^0+D^Tq=0$ and all edge complementarity conditions. The convex KKT conditions show that $z$ is the unique local constant-target projection. Consequently $2w^Ts^0-\norm w^2-2h_{\K_I}(w)=\norm w^2$, and direct summation proves \eqref{eq:Lc}.

If the block begins at slot $x+1$, its actual target is
\begin{equation}
 s_i=\epsilon_i[a+\eta(t_i-c)],\qquad a=\rho_0+\eta(x+c).
\end{equation}
Both $w$ and $\epsilon$ are reflection-symmetric, while $t_i-c$ is antisymmetric. Hence
\begin{equation}
 \sum_iw_i\epsilon_i(t_i-c)=0,
 \quad w^Ts=w^Ts^0,
 \quad\norm s^2=2na^2+\eta^2M_2,
\end{equation}
where $M_2=\sum_{i\in I}(t_i-c)^2=O((n+k)^3)$. This is an exact affine cancellation; the same vector is a feasible dual certificate for the affine target.

Concatenate the local dual directions, assigning zero weight to the healthy prefix, any early blocks that fail the amplitude condition, and the final unused reads. Every feasible whole-history path restricts to a feasible path on each block. Therefore \eqref{eq:support} holds even if the particular local primal optimizers cannot be joined. The elementary weak dual inequality
\begin{equation}
 \min_{z\in\K}\norm{s-z}^2
 \ge2w^Ts-\norm w^2-2h_\K(w)
\end{equation}
then proves \eqref{eq:closedG} and
\begin{equation}
2\sigma^2D_H\le E_{\rm dead}+E_{\rm tail}
+\sum_\ell[L_c(a_\ell,n_\ell)+\eta^2M_{2,\ell}].\label{eq:sharpfinite}
\end{equation}
This concatenates dual directions and uses restriction of the global feasible class.

For the terminal-balanced construction, $M=\sqrt{H/\xi}+O(1)$, $n_\ell=\xi\ell+O(1)$, and $a_\ell=\eta\xi\ell^2+O(\ell+1)$. Each block is perturbed by $O(1)$ readings, and fewer than four tail slots remain. The leading sums are
\begin{align}
 E_{\rm dead}&=\frac{2k\eta^2}{5\sqrt\xi}H^{5/2}+O(H^2),\nonumber\\
 \sum_\ell L_c(a_\ell,n_\ell)&=\frac{r\eta\sqrt\xi}{5}H^{5/2}+O(H^2),\nonumber\\
 \eta^2\sum_\ell M_{2,\ell}+E_{\rm tail}&=O(H^2).
\end{align}
Indeed, the first two per-block main terms are $2k\eta^2\xi^2\ell^4$ and $r\eta\xi^3\ell^4$; all other terms sum to $O(M^4)$. The finite number of early zero-weight blocks changes only the constant. Substitution in \eqref{eq:sharpfinite} proves \eqref{eq:upper}. Its coefficient is minimized at $\xi=2k\eta/r$, where it equals $2\sqrt2\sqrt{kr}\eta^{3/2}/5$. Combining with \eqref{eq:sharp_lower} proves the exact limit.
\end{proof}

